% method.tex -- Stochastic Attention via Langevin Dynamics on the Hopfield Energy

Section~3 established that attention is gradient descent on the modern Hopfield energy $E$, that $E$ is smooth and confining, and that Langevin dynamics converts any such energy into a sampler. We now combine these facts to derive \emph{stochastic attention}.

\subsection{Stochastic Attention Update}
The memory matrix $\mathbf{X}=[\mathbf{m}_1,\dots,\mathbf{m}_K]\in\mathbb{R}^{d\times K}$ stores $K$ patterns. We sample from the Boltzmann distribution $p_\beta(\boldsymbol{\xi})\propto\exp(-\beta E(\boldsymbol{\xi}))$ by applying the ULA~\eqref{eq:ula} to $U=\beta E$. Using $\nabla E(\boldsymbol{\xi})=\boldsymbol{\xi}-\mathbf{T}(\boldsymbol{\xi})$ and substituting $\tilde\alpha=\alpha/\beta$ so that the effective gradient step size is $\alpha$, we obtain
\begin{equation}\label{eq:stochastic-attention}
    \boldsymbol{\xi}_{t+1}
    = \boldsymbol{\xi}_t
      - \alpha\bigl(\boldsymbol{\xi}_t - \mathbf{T}(\boldsymbol{\xi}_t)\bigr)
      + \sqrt{\tfrac{2\alpha}{\beta}}\;\boldsymbol{\epsilon}_t,
    \qquad \boldsymbol{\epsilon}_t\sim\mathcal{N}(\mathbf{0},\mathbf{I}),
\end{equation}
which, after collecting terms, takes the form
\begin{equation}\label{eq:stochastic-attention-expanded}
    \boxed{\;
    \boldsymbol{\xi}_{t+1}
    = (1-\alpha)\,\boldsymbol{\xi}_t
      + \alpha\,\mathbf{X}\,\operatorname{softmax}\!\bigl(\beta\,\mathbf{X}^{\top}\boldsymbol{\xi}_t\bigr)
      + \sqrt{\tfrac{2\alpha}{\beta}}\;\boldsymbol{\epsilon}_t
    \;}
\end{equation}
This is the \emph{stochastic attention} update: a contraction toward the origin, a softmax-weighted pull toward stored memories, and isotropic Gaussian noise governed by $1/\beta$. The complete procedure is given in Algorithm~\ref{alg:stochastic-attention}. Because $\mathbf{X}$ is fixed, each step requires only the same primitives as a single attention head (two matrix-vector products, one softmax, one Gaussian draw) at cost $\mathcal{O}(dK)$. No training loop, score network, or learned parameters are needed.

\begin{algorithm}[t]
\caption{Stochastic Attention Sampler}\label{alg:stochastic-attention}
\begin{algorithmic}[1]
\REQUIRE Memory matrix $\mathbf{X}\in\mathbb{R}^{d\times K}$, inverse temperature $\beta>0$, step size $\alpha\in(0,1)$, number of iterations $T$, initial state $\boldsymbol{\xi}_0\in\mathbb{R}^d$
\ENSURE Sample trajectory $\boldsymbol{\xi}_0,\boldsymbol{\xi}_1,\dots,\boldsymbol{\xi}_T$
\FOR{$t = 0,1,\dots,T-1$}
    \STATE $\mathbf{a}_t \leftarrow \operatorname{softmax}\!\bigl(\beta\,\mathbf{X}^{\top}\boldsymbol{\xi}_t\bigr)$ \hfill $\triangleright$ attention weights
    \STATE $\boldsymbol{\epsilon}_t \sim \mathcal{N}(\mathbf{0},\mathbf{I}_d)$
    \STATE $\boldsymbol{\xi}_{t+1} \leftarrow (1-\alpha)\,\boldsymbol{\xi}_t + \alpha\,\mathbf{X}\,\mathbf{a}_t + \sqrt{2\alpha/\beta}\;\boldsymbol{\epsilon}_t$
\ENDFOR
\end{algorithmic}
\end{algorithm}

\subsection{Properties and Limiting Behavior}
When $\beta\to\infty$ and noise vanishes, the softmax sharpens to $\arg\max$ and the update reduces to deterministic retrieval (\emph{hard attention}). At finite $\beta$ with zero noise, it yields the standard softmax-weighted average (\emph{soft attention}). With finite $\beta$ and noise, the sampler generates novel patterns shaped by the memory geometry (\emph{stochastic attention}). The step size $\alpha\in(0,1)$ controls discretization fidelity versus mixing speed.

The score function $\nabla\log p_\beta = \beta(\mathbf{T}-\boldsymbol{\xi})$ is known in closed form from a single attention operation; no auxiliary network is needed. The following proposition (proved in Appendix~\ref{app:proof-regularity}) makes the regularity conditions explicit.

\begin{proposition}\label{prop:regularity}
Let $\sigma_{\max}=\|\mathbf{X}\|_{\mathrm{op}}$ denote the largest singular value of the memory matrix. The modern Hopfield energy~\eqref{eq:modern-hopfield-energy} has Lipschitz-continuous gradient with constant $L = 1 + \beta\sigma_{\max}^{2}/2$ and satisfies the dissipativity condition $\langle \nabla E(\boldsymbol{\xi}),\boldsymbol{\xi}\rangle \ge \|\boldsymbol{\xi}\|_2^2/2 - M^2/2$ for all $\boldsymbol{\xi}\in\mathbb{R}^d$.
\end{proposition}

The Hessian satisfies $\nabla^2 E \succeq (1-\beta\sigma_{\max}^{2}/2)\mathbf{I}$, so $E$ is strictly convex when $\beta\sigma_{\max}^{2}<2$, yielding a standard convergence result:

\begin{corollary}\label{cor:convergence}
When $\beta\sigma_{\max}^{2}<2$, the potential $U=\beta E$ is $m$-strongly convex with $m=\beta(1-\beta\sigma_{\max}^{2}/2)$ and has $L_{U}$-Lipschitz gradient with $L_{U}=\beta(1+\beta\sigma_{\max}^{2}/2)$. The iterates of Algorithm~\ref{alg:stochastic-attention} converge to~$p_\beta$ in $W_2$ at a geometric rate, with discretization bias vanishing as $\alpha\to 0$~\cite{dalalyanTheoreticalGuaranteesApproximate2017,durmusNonasymptoticAnalysisUnadjusted2017}.
\end{corollary}

In our experiments $\sigma_{\max}^{2}\approx 2$--$3$, so the convex regime holds only for $\beta\lesssim 1$, well below our operating range ($\beta\in\{200,2000\}$). We state this boundary explicitly because it delineates where the Hopfield energy transitions from unimodal (one global basin) to multimodal (one local minimum per stored pattern): the condition $\beta\sigma_{\max}^{2}=2$ is the onset of non-convexity, beyond which energy barriers separate stored memories.

For general non-convex energies, mixing time scales exponentially in the barrier height divided by the temperature~\cite{bovier2004metastability}. In the Hopfield energy at high~$\beta$, barriers between patterns are $O(\beta)$ while the temperature is $1/\beta$, so inter-basin mixing is exponentially slow. This is consistent with our empirical observation that a single chain at $\beta{=}2000$ does not cross between basins (single-chain diversity $0.282$; Appendix~\ref{app:single-chain}). However, within each basin the energy is locally strongly convex (the Hessian near each minimum is positive definite), so \emph{intra-basin} mixing is fast. The multi-chain protocol (Section~5) exploits this: by initializing chains near distinct stored patterns, we obtain cross-basin coverage through initialization and rely on fast local mixing within each basin. At $\beta{=}200$, the barriers are low enough that single chains do cross basins (single-chain diversity $0.796$, exceeding the multi-chain $\beta{=}2000$ value of $0.600$; Appendix~\ref{app:single-chain}).

The confining bound~\eqref{eq:energy-lower-bound} keeps iterates concentrated near $\mathrm{conv}\{\mathbf{m}_1,\dots,\mathbf{m}_K\}$ with excursions controlled by $\sqrt{2\alpha/\beta}$. Two guarantees hold for all $\beta>0$ regardless of convexity: the continuous-time SDE is geometrically ergodic by dissipativity~\cite{robertsTweedieExponentialConvergence1996}, and the MALA acceptance rate of $99.2\%$ at $\alpha{=}0.01$ (Table~\ref{tab:mnist-baselines}) confirms that ULA discretization bias is negligible in practice.

\subsection{Phase Transition and Temperature Selection}\label{sec:phase-transition}

The attention entropy $H(\beta) = -\sum_k p_k \log p_k$, where $p_k = \operatorname{softmax}(\beta\mathbf{X}^{\top}\boldsymbol{\xi})_k$, quantifies selectivity for a given state~$\boldsymbol{\xi}$. At $\beta=0$ the distribution is uniform ($H=\log K$); as $\beta\to\infty$ it concentrates on the nearest memory ($H\to 0$).

\begin{proposition}\label{prop:entropy-derivative}
For fixed $\boldsymbol{\xi}\in\mathbb{R}^d$, let $e_k = \mathbf{m}_k^{\top}\boldsymbol{\xi}$ and $\mathbf{p}=\operatorname{softmax}(\beta\mathbf{e})$. Then
\begin{equation}\label{eq:entropy-deriv}
\frac{dH}{d\beta} \;=\; -\beta\,\mathrm{Var}_{\mathbf{p}}(e),
\end{equation}
where $\mathrm{Var}_{\mathbf{p}}(e) = \sum_k p_k\,e_k^2 - \bigl(\sum_k p_k\,e_k\bigr)^{\!2}$.
\end{proposition}

\begin{proof}
Write $\log p_k = \beta e_k - \log Z$ with $Z=\sum_j\exp(\beta e_j)$, so $H = -\beta\langle e\rangle_{\mathbf{p}} + \log Z$. Differentiating and using $d(\log Z)/d\beta = \langle e\rangle_{\mathbf{p}}$ and $d\langle e\rangle_{\mathbf{p}}/d\beta = \mathrm{Var}_{\mathbf{p}}(e)$ yields the result.
\end{proof}

Since $\mathrm{Var}_{\mathbf{p}}(e)\ge 0$, entropy decreases monotonically. The derivative vanishes at both $\beta=0$ (prefactor) and $\beta\to\infty$ (variance collapse), so the maximum rate of entropy loss occurs at an inflection point $\beta^{*}$ satisfying $d^{2}H/d\beta^{2}=0$:
\begin{equation}\label{eq:inflection}
\mathrm{Var}_{\mathbf{p}}(e)\big|_{\beta^{*}} \;=\; -\beta^{*}\;\mu_{3}\big|_{\beta^{*}},
\end{equation}
where $\mu_3 = \langle(e - \langle e\rangle_{\mathbf{p}})^3\rangle_{\mathbf{p}}$ is the third central moment. This inflection marks the retrieval--generation phase boundary: below $\beta^{*}$, attention remains diffuse (generation); above it, attention concentrates on individual memories (retrieval).

For random memories on $\mathbb{S}^{d-1}$, the similarities $e_k$ have $\mathrm{Var}(e_k)=1/d$, so the softmax logits have standard deviation $\beta/\sqrt{d}$; the transition occurs when this reaches order unity, giving $\beta^{*}\sim\sqrt{d}$. Substituting into the per-step signal-to-noise ratio (Eq.~\ref{eq:snr}) yields
\begin{equation}\label{eq:snr-star}
\mathrm{SNR}^{*} \;=\; \sqrt{\frac{\alpha\,\beta^{*}}{2d}} \;\sim\; \sqrt{\frac{\alpha}{2\sqrt{d}}}\,.
\end{equation}
At $d{=}64$ and $\alpha{=}0.01$ this evaluates to $\mathrm{SNR}^{*}=0.025$, matching the empirical transition in Section~5. For structured data the effective variance may differ from $1/d$, but Eq.~\eqref{eq:inflection} provides a fully data-dependent criterion: evaluate $d^{2}H/d\beta^{2}$ over a $\beta$ sweep and locate the zero crossing.
